% Висновки залежать від дослідженого архіву, міток і меж вимірювання.
\section{Висновки}\label{sec:conclusion}

Дослідження робить пошукові рішення archive-index доступними для перевірки завдяки побудованій багатомовній колекції, замороженим запитам і прогонам, збереженим модельним оцінкам та генерованим звітам. За основними переглянутими мітками Gemini на 170 запитах із відповіддю середнє nDCG@10 становить 0,674 для E5, 0,633 для EG2, 0,629 для адаптивного злиття, 0,482 для BM25 із лемами та 0,450 для звичайного BM25. Адаптивне злиття перевищує BM25 після поправки, проте ці дані не встановлюють переваги над EG2 чи RRF. Свідчення користі морфології залежать від судді, а окреме розширення показує зміну точкового порядку енкодерів між наборами.

Ширший результат --- виміряний опис компромісів реалізації й оцінювання. Пошук із квантованими збірками залежить від охоплення оцінками та відмінностей входів/бекендів; синтетичне масштабування відділяє оптимізований відбір векторів від фактичного шляху ранжування архіву; мультимодальні завдання та перевірки приватного архіву відповідають на вужчі питання, ніж ефективність текстового пошуку для людини. Результати стосуються побудованого бенчмарку з модельними оцінками релевантності. Незалежні людські оцінки, незалежно сформульовані інформаційні потреби й перевірка на різних архівах залишаються майбутньою валідацією перед загальними твердженнями про пошук в особистих архівах.
